\scanpage{13}
\mspage{6}
On va introduire
\[
(M6)\quad\deleted{\unsure{\exists f_\alpha}}\quad
\left\{\begin{aligned}
&\forall \deleted{\unsure{X}}\inserted{Z},X\in\mathcal M_n,\quad Z\subset X,\quad\exists(f_\alpha)_{\alpha\in A}\ \text{famille}\\
&\text{finie de fonctions modérées }X\longrightarrow R\\
&\text{et des }(e_\alpha)_{\alpha\in A},\quad e_\alpha\in R_0,\quad\text{telle que}\\
&Z=\bigcap f_\alpha^{-1}(e_\alpha).\quad\text{De plus, OPS que les}\\
&f_\alpha\text{ séparent les points de }X\smallsetminus Z.
\end{aligned}\right.
\]

\margin{\uline{Prop} $(M6)$, $(M7)$ impliquent l'existence de la somme amalg. avec une flèche libre.}
\margin{\diagram{\begin{tikzcd}[column sep=small]
X_1&Z\arrow[l,hook']\arrow[r]&X_2
\end{tikzcd}}}

\deleted{Cela implique que l'on peut faire un diagramme \inserted{cart}}
\deleted{\diagram{\begin{tikzcd}
\deleted{X}\arrow[r,hook]&\deleted{R^n\times R^A}\\
\deleted{Z}\arrow[u,hook]\arrow[r,hook,no head]&\deleted{R^n\times\{e\}}\arrow[u,hook]
\end{tikzcd}}}
\deleted{$e=(e_\alpha)_{\alpha\in A}$}

\deleted{Appliquons ceci : $Z_1\hookrightarrow X_1$ et $Z_2\hookrightarrow X_2$}

ce qui \uline{équivaut} au fait que le quotient $X/Z$ (contraction de $Z$ en un point) existe dans $\mathcal E$ au sens strict (i.e. $X\mathrel{\mapstochar\relbar}|X|$ y commute). De plus, on suppose
\margin{\unsure{supplém.} important)}
\[
(M7)\quad\forall Z,X\in\mathcal M_n,\quad Z\subset X,\quad\text{et tte}
\]
\[
f:Z\longrightarrow R\text{ modérée},\qquad\exists g:X\longrightarrow R\text{ modérée}
\]
qui prolonge \unsure{$f$}

\uline{Proposition} Moyennant $M5,M6,M7$, les sommes amalgamées $X_1\amalg_Z X_2$ ($Z\hookrightarrow X_1$, $Z\hookrightarrow X_2$) existent dans $\operatorname{Esp}_{\mathcal M_*}$, et commutent au foncteur oubli\ldots
\margin{cf généralisation qui vient ci-dessous.}
\mspage{12}
